\section{Allgemeines zur Bibel}

Am letzten Donnerstag, habe ich hier, den Vortrag: \enquote{Ist die Bibel ein Märchenbuch?} gehalten. In diesem Vortrag ging es darum, wie die Bibel enstanden ist und was dieses Buch so einzigartig macht.

Diese Serie jetzt über die Bibel, geht in eine andere Richtung. Hier gehe ich davon aus, dass alle, die mir zuhören glauben, dass der Inhalt in diesem Buch von Gott inspiriert ist und der absoluten Wahrheit entspricht. Sieht das jemand nicht so, so kann er ruhig zuhören, aber keine Disskusion wird mich von meiner Meinung abbringen.

Wieso? Ich komme aus dem katholischen. Ich fühlte mich dort verarscht, weil das, was sie sagten und das was sie machten, nicht übereinstimmte. Auf meinem Lebensweg, kamen mir viele atheistischen Schriften in die Hände. Auch da gab es so viele Widersprüche. Ich fand einfach keine Wahrheit. War ja eigentlich auch egal, weil einen eh die Würmer fressen werden.

Nachdem ich Christus angenommen habe, fing ich an, geleitet durch den Heiligen Geist, die Bibel zu lesen. Wieso sage ich geleitet durch den Heiligen Geist? Weil ich die Bibel schon vor meiner Bekehrung gelesen habe, aber nichts verstand und nur Mord und Todschlag darin fand.

Nach meiner Bekehrung war das anders. Es wurde eine Neugierde in mir geweckt. Ich stellte Fragen, an Janina an den Pastor in der FEG und suchte im Internet.

Dieses Buch ist die einzige Quelle, die ich und wir besitzen, um die Wahrheit Gottes zu erfahren. 

Nochmals das Zitat von John McArthur:
\begin{quote}
    \textit{Die Lehre von der Schrift ist absolut fundamental und essenziell, da sie die einzige wahrhaftige Quelle aller christlichen Wahrheit ist.}
\end{quote}
Darauf will ich meinen Lehrteil die nächsten Sonntage aufbauen.
Aus diesen Fragen gibt es folgendes Inhaltsverzeichnis:
\begin{enumerate}
    \item Einführung in die Schrift
    \item Die inspiration der Schrift
    \item Die Autorität der Schrift
    \item Die Irrtumslosigkeit der Schrift
    \item Die Bewahrung der Schrift
    \item Das Lehren und Predigen der Schrift
    \item Die Verpflichtung gegenüber der Schrift
\end{enumerate}


\begin{otsblock}[Einführung in die Schrift]
    Woher weiss ich das der Inhalt in diesem Buch von Gott höchst persönlich ist? Hat mir das der Mitternachtsruf gesagt, oder Roger Liebi oder ein schlaues Theologiebuch?

    Nein, das hat mir Gott gesagt. Und nein, mir ist kein Engel erschienen und hat mir das erzählt, sondern es steht in der Bibel.

    Es wird in der Schrift wiederholt behauptet, dass sie Gottes Wort ist. Für die Propheten war die Schrift das Fundament, für Christus und sein Apostel gründete die gesamtheit der christlichen Lehrer auf dieser Schrift.\\
    Im alten Testament wird über 2500-mal gesagt, dass Gott gesprochen hat. Im Neuen Testament kommt der Ausdruck \betonung[b, p]{Gottes Wort} über 40-mal vor.

    Das Wort Gottes ist ein Schwerpunkt in der Apostelgeschichte und Lukas schreibt, wie sich das Wort schnell ausbreitete \biblezit{Apg}{6:7}, \biblezit{Apg}{12:24}, \biblezit{Apg}{19:20}.

    Die Bibel beansprucht höchste geistliche Autorität in Lehre, Überführung, Zurechweisung und Unterweisung in der Gerechtigkeit, weil sie das inspirierte Wort des allmächtigen Gottes ist. \biblezit{IITim}{3:16-17}

    Gottes Wort selber erklärt, dass es irrtumslos und unfehlbar ist.

    Wie wir in den letzten Sonntagen oft gehört haben, aber wir nicht oft genug hören können, ist, dass Gott wahrhaftig, fehlerlos, und vertrauenswürdig ist. Dasselbe gilt für sein Wort. Was jemand über Gottes Wort denkt, spiegelt wider, was er über Gott denkt.

    Es gibt sieben wesentliche Wesenszüge der Schrift:
    \begin{enumerate}
        \item wirksam \biblezit{IThes}{2:13}; \biblezit{Heb}{4:12}
        \item zuverlässig \biblezit{Jes}{55:10-11}; \biblezit{Lk}{16:17}
        \item kraftvoll \biblezit{Rom}{1:16-17}; \biblezit{IKor}{1:18}
        \item lebendig \biblezit{Joh}{6:63}; \biblezit{Heb}{4:12}; \biblezit{IPetr}{1:23}
        \item reinigend \biblezit{Eph}{5:26}
        \item nährend \biblezit{1Petr}{2:2}
        \item heiligend \biblezit{Joh}{17:17-19}
    \end{enumerate}
\end{otsblock}
\newpage